\section{工程案例：Agent 参与的真实仓库}
\subsection{半自动化提交}
Graham 分享了自己的仓库实践：有些项目中“一半的提交内容都是先由编码 Agent 草拟，再由他去检查与修复”的。Agent 先读 Issue，通过摘要与检索找到相关文件，然后草稿出初步补丁，最后由人类工程师合并。

他预计未来一两年内，Agent 会承担更多日常琐碎任务，让工程师把注意力放在高层设计。现在 Agent 还能顺利处理共享依赖的 Package 脚本、基础 CLI 操作与简单的 bug fix，复杂场景仍然需要人工复审。

\subsection{过程回顾与反馈}
每次 Agent 运行后，团队都会记录最终日志：包括主动 fetch 的文件、调用的命令、测试结果与失败原因。这些日志被转成训练样本与回调提醒，让 Agent 学会把失败样本中有用的步骤嵌入提示链。

他还强调：“观察 Agent 运行的过程并给出即时纠正，比在结果上打分更能提升模型质量。”这种观察式反馈包括人工暂停、注入提示、在纠错后再让 Agent 重试。

\begin{knowledgebox}{案例经验}
Graham 已经在多个仓库里使用 Agent 提笔，用 Agent 贡献的提交占据约一半工作量。这证明 Agent 在重复性任务中已具备实用价值。
\end{knowledgebox}

\begin{importantbox}{实用建议}
先把 Agent 投放在 10 分钟以内的任务（如格式化、补丁、依赖升级），拼凑人类 review 的流程，再逐渐向 longer-running job 扩展。
\end{importantbox}

\begin{warningbox}{不要让 Agent 独自部署}
即便 Agent 可以完成某些改动，也不应让其单独推到生产，必须保留人工审查、测试回归与手动 approve。
\end{warningbox}

\subsection{本章小结}
\begin{itemize}
    \item Agent 在真实仓库中的尝试已展现半自动化的潜力，但还不足以完全替代人类工程师。
    \item 记录过程日志与即时纠错是提升 Agent 可靠性的关键做法。
\end{itemize}

